\section{Sharp threshold and rational degrees}
Use signs \(\{-1,1\}\) for the raw table entries. The threshold degree of the promise is the least degree of a polynomial having opposite strict signs on the two promised sets. For rational approximation we require a quotient with positive denominator and values in \([0,1]\) on the entire Boolean cube, with error at most \(1/3\) on the promise. Its degree is the maximum of the two polynomial degrees. Allowing a nonvanishing denominator of arbitrary sign changes the sign-degree lower comparison by at most a factor of two, on multiplying by that denominator.

\begin{theorem}
For \(s=\Theta(N^a),\ t=\Theta(N^b)\) with fixed \(0<a<b<1\), the threshold degree and the globally bounded rational approximate degree of gap MCSP are both \(\Theta_{a,b}(s)\).
\end{theorem}
\begin{proof}
Write \(M=|A_s|\). Apply the preceding minimax argument to \(A_s\cup(-A_s)\subseteq A_{s+1}\), accounting for a possibly charged output negation. This gives, for each hard table \(h\), a distribution \(\mu_h\) on positions with
\[
  \left|\sum_i\mu_h(i)h_i y_i\right|\le\epsilon
       \quad(y\in A_s),\qquad \epsilon=O(\sqrt{(s+1)n/t}).
\]
A disagreement lower bound also bounds agreement from below by comparison with the complemented circuit. Choose the constant in \(\epsilon\) large enough; for large \(n\) it is less than one. Take
\[
 d=\left\lceil\frac{\log(20M)}{2\log(1/\epsilon)}\right\rceil=O_{a,b}(s).
\]
The bound uses \(\log M=O(s\log s)\) and \(\log(1/\epsilon)=\Omega_{a,b}(\log N)\). Let \(\phi(x)\) be the vector of all squarefree parity monomials through degree \(d\), including the constant, and put
\[
  K=M^{-1}\sum_{y\in A_s}\phi(y)\phi(y)^{\mathsf T},\quad
  G=K+\epsilon^{2d}I,\quad
  P(x)=\phi(x)^{\mathsf T}G^{-1}\phi(x).
\]
All matrices are finite and \(G\) is positive definite. The polynomial \(P\) is nonnegative and has degree at most \(2d\). Each easy table \(y\) has \(P(y)\le M\), since \(G\succeq\phi(y)\phi(y)^{\mathsf T}/M\); conjugate by \(G^{-1/2}\).

For hard \(h\), expand
\[
     q_h(x)=\left(\sum_i\mu_h(i)h_ix_i\right)^d
            =v_h^{\mathsf T}\phi(x)
\]
on the cube. Before reduction the sum of absolute multinomial coefficients is \(1\). Using \(x_i^2=1\) and collecting terms gives \(\|v_h\|_2\le\|v_h\|_1\le1\).
Therefore \(v_h^{\mathsf T}Gv_h\le 2\epsilon^{2d}\), since \(q_h^2\le\epsilon^{2d}\) on every easy table. But \(q_h(h)=1\). Matrix Cauchy--Schwarz yields
\(P(h)\ge 1/(2\epsilon^{2d})\ge10M\).
Thus \(P-3M\) strictly separates the two classes, and \(3M/(P+3M)\) is a globally bounded rational approximant with the required orientation and accuracy.

For the lower bound use the one-mark construction above, with \(k=\Theta(t)\) buckets of size \(H=N/k\). The tuples of easy labels form a binary linear subgroup \(C\subseteq(\mathbb F_2^{\log_2H})^k\), since polynomial evaluation and coordinate projection are linear in their coefficients. Its logarithmic cardinality is \(O(sn)\) bits. Translate it by a uniformly random label tuple \(v\). Each fixed element of \(v+C\) is then a uniform label tuple, and a uniform tuple's marked table has size at most \(t\) with probability at most \(|A_t|/H^k\). The constant in \(k/t\) was chosen with fixed slack, so \(|C||A_t|/H^k=o(1)\). Some coset therefore consists entirely of hard tables. Both uniform laws on \(C\) and on this coset give independent uniform labels in every set of at most \(D=\Theta(s)\) buckets. Every raw-table monomial of degree at most \(D\) has identical expectations on the two laws, because it touches at most \(D\) buckets. A strictly sign-separating polynomial of that degree would have opposite expected signs. Finally, subtracting \(1/2\) from a rational approximant and clearing its positive denominator gives a sign polynomial. This proves both lower bounds.
\end{proof}
The positive quadratic form \(P\) can take large values away from the easy set. The rational denominator supplies the global boundedness used here; the proof gives a distinct upper bound for a different approximation model than bounded ordinary polynomials.

The bounded-error quantum query complexity with postselection is consequently
$\Theta(s)$. For the upper bound apply the amplitude construction of
Mahadev and de Wolf to the numerator and denominator just obtained.
Its success analysis is pointwise, so the approximation guarantee is needed
only on the promise, and the denominator is positive everywhere.
Conversely, if a $q$-query algorithm postselects with positive probability
on its promise, its joint acceptance and postselection probability $A(x)$
and its postselection probability $B(x)$ are polynomials of degree at most
$2q$. The polynomial $2A-B$ strictly separates the two promise classes,
which gives the lower bound even when postselection can fail identically
off promise. See U.~Mahadev and R.~de Wolf,
\emph{Rational approximations and quantum algorithms with postselection},
Quantum Information and Computation 15 (2015), 295--307.
